%% abstract.tex -- submission version

\begin{abstract}
Self-healing mechanisms capture execution failures and feed diagnostic content into a subsequent large-language-model (LLM) turn. We study the resulting recovery-channel prompt-injection path, in which externally influenced failure content is re-injected independently of the normal forward path. We combine systematic multiple-case source-code analysis with qualitative coding and controlled native-framework experiments.

Across five model aliases, the canonical Reflexion matrix reports attack success rates from \ReflexionMinASR{} to \ReflexionMaxASR{} with zero classified attacks in \ReflexionBenignTrials{} matched benign trials. Native AutoGen and Aider recovery/error loops provide additional framework evidence. In a native Reflexion port, raw recovery succeeds in \NativeTrustOriginRawSuccesses{}/\NativeTrustOriginAttackTrials{} trials (\NativeTrustOriginRawASR{}), compared with \NativeTrustOriginTaggedSuccesses{}/\NativeTrustOriginAttackTrials{} (\NativeTrustOriginTaggedASR{}) under TrustOrigin framing; both matched benign arms have zero false positives. This native contrast is the primary provenance result.

Supplementary screens report \DefenseUndefendedASR{} versus \DefenseDefendedASR{} attack success ($p=\DefenseFisherP{}$, $n=100$ per pooled configuration) and \LegalRecoveryRate{} legitimate-recovery success over \LegalRecoveryTasks{} tasks. A fixed local \QwenLocalModel{} sensitivity run contributes 100 valid trials (raw ASR \QwenLocalRawASR{}, TrustOrigin ASR \QwenLocalTaggedASR{}, zero benign false positives) and is explicitly exploratory. Hosted provider aliases expose no immutable checkpoint, so exact snapshot reproducibility remains a limitation.
\end{abstract}
